\documentclass[12pt]{article}

% -------------------------
% PAQUETES NECESARIOS
% -------------------------
\usepackage[utf8]{inputenc}
\usepackage[T1]{fontenc}
\usepackage{titlesec}
\usepackage{fancyhdr}
\usepackage{graphicx}
\usepackage{listings}
\usepackage{xcolor}
\usepackage{geometry}
\usepackage[spanish]{babel}
\usepackage{hyperref}

% -------------------------
% CONFIGURACIÓN DE PÁGINA
% -------------------------
\geometry{margin=2.5cm}

% Fuente tipo Times New Roman equivalente
\usepackage{mathptmx}

% Interlineado sencillo (por defecto en LaTeX)

% -------------------------
% CONFIGURACIÓN DE ENCABEZADO Y PIE
% -------------------------
\pagestyle{fancy}
\fancyhf{}

\fancyhead[L]{Multiprogramación paralela}
\fancyhead[R]{2026/2027}

\fancyfoot[L]{Andrey, Abdellah} % <-- EDITAR
\fancyfoot[R]{\thepage}

\renewcommand{\headrulewidth}{0.4pt}
\renewcommand{\footrulewidth}{0.4pt}

% -------------------------
% INICIO DEL DOCUMENTO
% -------------------------
\begin{document}

% =========================
% PORTADA
% =========================
\thispagestyle{empty}

% Bloque arriba a la derecha
\begin{flushright}
Abdellah Bensaidi, Souri \\
Morales Vicuña, Andrey \\
Asignatura: Multiprogramación paralela \\
Grupo: 1 \\
%  Subgrupo: 1.1 \\

\end{flushright}

\vfill

% Imagen centrada
\begin{center}
\includegraphics[width=0.4\textwidth]{sello.png}
\end{center}

\vfill

% Fecha abajo
\begin{center}
\today
\end{center}

\newpage

\newpage

\begin{enumerate}
    \item \textbf{}

    \item \textbf{}

    \item \textbf{}

    %\begin{itemize}
    %    \item \textbf{}
    %    \item \textbf{}
    %\end{itemize}

    %\begin{verbatim} 
    %\end{verbatim}


\end{enumerate}


\begin{thebibliography}{9}

\bibitem{baeldung_parallel}
Baeldung, \emph{A Quick Guide to the Linux parallel Command}, 2025. [En línea]. Disponible en: \url{https://www.baeldung.com/linux/parallel-command-tutorial} [Accedido: 09-sep-2026].

\bibitem{lemon_tendencias}
Lemon Learning, \emph{10 tendencias IT}, s.f. [En línea]. Disponible en: \url{https://lemonlearning.com/es/blog/10-tendencias-it} [Accedido: 09-sep-2026].

\end{thebibliography}










\end{document}